\section{Aplicaciones y limitaciones de la IA en las pruebas de seguridad}

\subsection{Formas en que la IA mejora la seguridad}

\begin{itemize}
    \item El «shift left» de la seguridad es más fácil de lograr que antes
    \item Los LLM pueden detectar problemas dentro del flujo de trabajo
    \item Pruebas de penetración automatizadas (automated penetration testing)
\end{itemize}

\subsection{Limitaciones graves actuales}

\begin{warningbox}{El problema de la tasa de falsos positivos en el AI SAST}
En el AI SAST, la tasa de falsos positivos es muy alta: la de Claude Code/Codex puede llegar a 50--100\% (según el tipo de vulnerabilidad), y la de las técnicas de SAST tradicionales también supera el 50\%. Esto significa que una gran cantidad de alertas de seguridad son falsas alarmas, lo cual interfiere seriamente con el flujo de trabajo de los desarrolladores.
\end{warningbox}

\begin{itemize}
    \item \textbf{Benchmarks poco realistas}: los benchmarks existentes con frecuencia no reflejan escenarios reales, lo que dificulta evaluar la capacidad de análisis de seguridad de los LLM
    \item \textbf{Análisis no determinista}: ejecutar el mismo prompt varias veces produce resultados distintos; ¿cómo se garantiza capturar todas las vulnerabilidades?
    \item \textbf{Context rot}: no todo el context tiene la misma utilidad
    \item \textbf{Pérdida de información por compaction}: la compresión mediante resúmenes para ajustarse a la context window puede perder detalles de seguridad clave
\end{itemize}

\subsection{Preguntas abiertas}

\begin{enumerate}
    \item ¿Cómo reducir los falsos positivos y las hallucinations en la detección de vulnerabilidades?
    \item ¿Cómo verificar que un parche generado por un LLM es seguro y no introduce regresiones?
    \item ¿Puede un LLM explicar \textbf{por qué} marcó cierta vulnerabilidad o propuso cierta solución?
    \item ¿Cuál es el benchmark correcto para medir el desempeño de un LLM en AppSec?
    \item ¿Cómo debería integrarse un LLM en el flujo de CI/CD sin inundar al equipo de ruido?
    \item Si un parche generado por IA introduce una vulnerabilidad, ¿quién debería responder por ello?
\end{enumerate}

\subsection{Integrar las verificaciones de seguridad de IA en el flujo de desarrollo}

Al llevar el contenido del curso a la práctica de ingeniería, lo más importante no es agregar otro «escáner inteligente», sino colocarlo en el lugar adecuado:

\begin{table}[H]
\centering
\begin{tabular}{p{0.18\textwidth} p{0.24\textwidth} p{0.34\textwidth}}
\toprule
Etapa & Acción de seguridad de IA adecuada & Objetivo principal \\
\midrule
Codificación & Sugerencias en línea dentro del IDE, verificaciones estáticas locales & Impedir lo antes posible que entren al repositorio patrones evidentemente peligrosos \\
Confirmación (commit) & Revisión de vulnerabilidades a nivel de PR, escaneo de cambios en dependencias & Detener cambios de alto riesgo antes de fusionarlos \\
Pruebas & DAST automatizado / apoyo de fuzzing & Descubrir rutas en tiempo de ejecución y problemas relacionados con el entorno \\
Después del lanzamiento & Resumen de alertas de monitoreo, apoyo en el análisis posterior a incidentes & Acortar el tiempo de detección y localización, y formar un ciclo de aprendizaje \\
\bottomrule
\end{tabular}
\caption{Puntos de integración de las capacidades de seguridad de IA en el flujo de desarrollo}
\end{table}

\begin{knowledgebox}{La estrategia más eficaz por lo general no es «dejar que la IA se haga cargo de la seguridad de forma independiente»}
Un enfoque más realista es colocar a la IA en una posición de «alta cobertura, baja capacidad de decisión»: su función es exponer riesgos lo antes posible, mientras que las personas y los sistemas de prueba rigurosos toman la decisión final. Así se aprovecha la cobertura de la IA y a la vez se controla el costo de ruido que generan los falsos positivos y las hallucinations.
\end{knowledgebox}

\subsection{Resumen de la sección}

La IA en el ámbito de las pruebas de seguridad representa tanto una oportunidad como un reto. Los LLM pueden reducir la dificultad de lograr el «shift left», pero todavía quedan por resolver la alta tasa de falsos positivos, la salida no determinista y los problemas de gestión del context. La atribución de responsabilidad es una cuestión ética clave que aún no tiene respuesta.

